\chapter{Introduction}
\label{ch:introduction}

\section{Motivation}
\label{sec:motivation}

Deep learning has transformed machine perception.
Convolutional neural networks now classify images, transcribe speech, and recognise environmental sounds at or above human-level accuracy on standard benchmarks~\citep{kong2020panns, gong2021ast}.
Yet this capability carries a substantial energy cost.
A single inference through a moderately sized convolutional network on a GPU consumes tens to hundreds of millijoules~\citep{yik2025neurobench}, while the human brain performs real-time, continuous auditory classification within an overall power budget of approximately 20\,W---sustaining on the order of $10^{16}$ synaptic operations per second~\citep{attwell2001energy}.
The disparity between biological and artificial efficiency spans three to four orders of magnitude for equivalent perceptual tasks.

This matters because the most compelling applications for audio intelligence are energy-constrained.
Hearing aids operate on roughly 1\,mW with a weekly charge cycle.
Wildlife monitoring sensors must persist for months on a single battery cell.
Smart building occupancy detectors, industrial fault listeners, and search-and-rescue robots all require continuous audio classification at the edge, where cloud offloading is impractical due to latency, bandwidth, and privacy constraints.
There exists, therefore, a genuine engineering gap: the algorithms that achieve state-of-the-art accuracy are far too expensive to run locally on the devices that need local inference most.

